% Part: methods
% Chapter: induction
% Section: relations

\documentclass[../../../include/open-logic-section]{subfiles}

\begin{document}

\olfileid[pt]{mth}{ind}{rel}

\olsection{Relações e Funções}

Quando definimos um conjunto de objetos (como os números naturais ou
os termos bonitos) indutivamente, também podemos definir \emph{relações sobre}
esses objetos por indução.  Por exemplo, considere a seguinte
ideia: um termo bonito~$t_1$ é um subtermo de um termo bonito~$t_2$ se ocorre como
parte dele.  Vamos usar um símbolo para isso: $t_1 \sqsubseteq t_2$. Todo
termo bonito é, é claro, um subtermo de si mesmo: $t \sqsubseteq t$. Podemos
dar uma definição indutiva dessa relação como segue:

\begin{defn}
  A relação de um termo bonito~$t_1$ ser um subtermo de~$t_2$, $t_1
  \sqsubseteq t_2$, é definida por indução em~$t_2$ como segue:
  \begin{enumerate}
  \item Se $t_2$ é uma letra, então $t_1 \sqsubseteq t_2$ sse $t_1 = t_2$.
  \item Se $t_2$ é $[s_1 \circ s_2]$, então $t_1 \sqsubseteq t_2$ sse $t_1 =
    t_2$, $t_1 \sqsubseteq s_1$, ou $t_1 \sqsubseteq s_2$.
  \end{enumerate}
\end{defn}

Esta definição, por exemplo, nos dirá que $\mathrm{a}
\sqsubseteq [\mathrm{b} \circ \mathrm{a}]$. Pois (2) diz que
$\mathrm{a} \sqsubseteq [\mathrm{b} \circ \mathrm{a}]$ sse $\mathrm{a}
= [\mathrm{b} \circ \mathrm{a}]$, ou $\mathrm{a} \sqsubseteq b$, ou
$\mathrm{a} \sqsubseteq \mathrm{a}$. As duas primeiras são falsas:
$\mathrm{a}$ claramente não é idêntico a $[\mathrm{b} \circ
  \mathrm{a}]$, e por~(1), $\mathrm{a} \sqsubseteq \mathrm{b}$ sse
$\mathrm{a} = \mathrm{b}$, o que também é falso. No entanto, também por~(1),
$\mathrm{a} \sqsubseteq \mathrm{a}$ sse $\mathrm{a} = \mathrm{a}$,
o que é verdadeiro.

É importante observar que o sucesso desta definição depende de
um fato que ainda não provamos: todo termo bonito~$t$ ou é uma
letra isolada, ou existem termos bonitos \emph{univocamente determinados}
$s_1$ e $s_2$ tais que $t = [s_1 \circ s_2]$.  ``Univocamente
determinados'' aqui significa que, se $t = [s_1 \circ s_2]$, ele não é
\emph{também} $= [r_1 \circ r_2]$ com $s_1 \neq r_1$ ou $s_2 \neq r_2$.
Se este fosse o caso, então a cláusula~(2) poderia entrar em conflito
consigo mesma: lendo $t_2$ como $[s_1 \circ s_2]$ poderíamos obter $t_1
\sqsubseteq t_2$, mas se lermos $t_2$ como $[r_1 \circ r_2]$ poderíamos
obter não $t_1 \sqsubseteq t_2$.  Antes de provarmos que isso não pode
acontecer, vejamos um exemplo em que ele \emph{pode} acontecer.

\begin{defn}
  Defina \emph{termos sem colchetes} indutivamente por
  \begin{enumerate}
  \item Toda letra é um termo sem colchetes.
    \item Se $s_1$ e $s_2$ são termos sem colchetes, então $s_1 \circ s_2$ é
      um termo sem colchetes.
    \item Nada mais é um termo sem colchetes.
  \end{enumerate}
\end{defn}

Termos sem colchetes são, por exemplo, $\mathrm{a}$, $\mathrm{b} \circ
\mathrm{d}$, $\mathrm{b} \circ \mathrm{a} \circ \mathrm{b}$. Agora, se
definíssemos ``subtermo'' para termos sem colchetes do modo como fizemos acima, a
segunda cláusula diria
\begin{center}
  Se $t_2 = s_1 \circ s_2$, então $t_1 \sqsubseteq t_2$ sse $t_1 = t_2$, $t_1
  \sqsubseteq s_1$, ou $t_1 \sqsubseteq s_2$.
\end{center}

Agora, $\mathrm{b} \circ \mathrm{a} \circ \mathrm{b}$ é da forma $s_1
\circ s_2$ com 
\begin{align*}
  s_1 & = \mathrm{b} \text{ e} & s_2 & = \mathrm{a} \circ \mathrm{b}.
\intertext{Ele também é da forma $r_1 \circ r_2$ com}
  r_1 & = \mathrm{b} \circ \mathrm{a} \text{ e} & r_2 &= \mathrm{b}.
\end{align*}
Agora, $\mathrm{a} \circ \mathrm{b}$ é um subtermo de $\mathrm{b} \circ
\mathrm{a} \circ \mathrm{b}$?  A resposta é sim, se seguirmos a primeira
leitura, e não, se seguirmos a segunda.

A propriedade de que o modo como um termo bonito é construído a partir de outros
termos bonitos é único é chamada de \emph{leitura única}. Como definições
indutivas de relações para tais objetos definidos indutivamente são
importantes, temos que provar que ela vale.

\begin{prop}
  Suponha que $t$ é um termo bonito. Então ou $t$ é uma letra isolada,
  ou existem termos bonitos univocamente determinados $s_1$, $s_2$ tais que~$t =
  [s_1 \circ s_2]$.
\end{prop}

\begin{proof}
  Se $t$ é uma letra isolada, a condição é satisfeita. Então suponha
  que $t$ não é uma letra isolada. Podemos perceber, pela definição
  indutiva, que então $t$ deve ser da forma $[s_1 \circ s_2]$ para alguns
  termos bonitos $s_1$ e~$s_2$. Resta mostrar que estes são univocamente
  determinados, isto é, se $t = [r_1 \circ r_2]$, então $s_1 = r_1$ e $s_2 = r_2$.

  Então suponha $t = [s_1 \circ s_2]$ e também $t = [r_1 \circ r_2]$ para termos bonitos
  $s_1$, $s_2$, $r_1$, $r_2$. Temos que mostrar que $s_1 = r_1$ e $s_2 = r_2$.
  Primeiro, $s_1$ e $r_1$ devem ser idênticos, pois, caso contrário, um é um
  segmento inicial próprio do outro. Mas, por \olref[sti]{prop:initial}, isso
  é impossível se $s_1$ e~$r_1$ são ambos termos bonitos.  Mas se $s_1 = r_1$,
  então claramente também $s_2 = r_2$.
\end{proof}

Também podemos definir funções indutivamente: por exemplo, podemos definir a
função~$f$ que mapeia qualquer termo bonito na profundidade máxima de aninhamento de
$[\dots]$ nele como segue:

\begin{defn}
  \ollabel{defn:depth} A \emph{profundidade} de um termo bonito, $f(t)$, é
  definida indutivamente como segue:
  \[
  f(t) = \begin{cases}
    0 & \text{ se $t$ é uma letra}\\
    \max(f(s_1), f(s_2)) + 1 & \text{ se $t = [s_1 \circ s_2]$.}
  \end{cases}
  \]
\end{defn}

Por exemplo,
\begin{align*}
  f([\mathrm{a} \circ \mathrm{b}]) & = 
    \max(f(\mathrm{a}),f(\mathrm{b})) + 1 = \\ 
  &= \max(0, 0) + 1 = 1, \text{ e}\\
  f([[\mathrm{a} \circ \mathrm{b}] \circ \mathrm{c}]) & = 
    \max(f([\mathrm{a} \circ \mathrm{b}]), f(\mathrm{c})) + 1 = \\ 
  & = \max(1,0) + 1 = 2.
\end{align*}

Aqui, é claro, supomos que $s_1$ e $s_2$ são termos bonitos, e fazemos
uso do fato de que todo termo bonito ou é uma letra ou é da forma
$[s_1 \circ s_2]$. É novamente importante que ele possa ser dessa forma de
apenas um modo. Para ver por quê, considere novamente os termos sem colchetes que
definimos antes. A ``definição'' correspondente seria:
\[
  g(t) = 
  \begin{cases}
    0 & \text{ se $t$ é uma letra}\\
   \max(g(s_1), g(s_2)) + 1 & \text{ se $t = s_1 \circ s_2$.}
  \end{cases}
\]
Agora considere o termo sem colchetes $\mathrm{a} \circ \mathrm{b} \circ
\mathrm{c} \circ \mathrm{d}$. Ele pode ser lido de mais de um modo,
por exemplo, como $s_1 \circ s_2$ com 
\begin{align*}
  s_1 & = \mathrm{a} \text{ e} & 
  s_2 & = \mathrm{b} \circ \mathrm{c} \circ \mathrm{d},
\intertext{ou como $r_1 \circ r_2$ com}
  r_1 & = \mathrm{a} \circ b \text{ e} & 
  r_2 &= \mathrm{c} \circ \mathrm{d}.
\end{align*}
Calcular $g$ de acordo com o primeiro modo de lê-lo daria
\begin{align*}
  g(s_1 \circ s_2) & =
  \max(g(\mathrm{a}), g(\mathrm{b} \circ \mathrm{c} \circ \mathrm{d})) + 1 =\\ & =
  \max(0,2) + 1 = 3
  \intertext{enquanto, de acordo com a outra leitura, obtemos}
  g(r_1 \circ r_2) & =
  \max(g(\mathrm{a} \circ \mathrm{b}), g(\mathrm{c} \circ \mathrm{d})) + 1=\\ &
  = \max(1,1) + 1 = 2
\end{align*}
Mas uma função deve sempre produzir um valor único; então nossa ``definição''
de~$g$ não define função alguma.

\begin{prob}
  Dê uma definição indutiva da função $l$, em que $l(t)$ é
  o número de símbolos no termo bonito~$t$.
\end{prob}

\begin{prob}
  Prove, por indução estrutural sobre termos bonitos $t$, que $f(t) < l(t)$ (em que
  $l(t)$ é o número de símbolos em~$t$ e $f(t)$ é a profundidade de
  $t$ tal como definida em \olref[mth][ind][rel]{defn:depth}).
\end{prob}

\end{document}
